% !TEX TS-program = lualatex+se+debug
% !TEX encoding = UTF-8 Unicode

\documentclass[14pt, twoside, a4paper]{extreport}
% -----------------------------------------------------------------------------------
% ------------------------- USTAWIENIA STRONY--------------------------
% -----------------------------------------------------------------------------------
\usepackage{geometry}
\geometry{
	left=10mm,
	right=10mm,
	top=15mm,
	bottom=15mm,
	marginparwidth=0mm,
	marginparsep=0mm
	}
\newcommand{\wstawplik}[1]{\input{../include/#1}}%skraca kod
\wstawplik{headerofficium}
% elementy zmienne co do headerofficium:
\graphicspath{{../img/}} % katalog obrazów
%\svgpath{{../include/}}
\grechangestaffsize{20}%rozmiar wyniku GregorioTEX
% koniec 
\gresetlastline{ragged}% justowanie do lewej
% debugowanie:
% %\iffalse %*** %\fi %***
% \tracingall przed problematycznym kodem - bardzo dużo informacji
% lualatex+se+debug - dodatkowe opcje: -file-line-error i -halt-on-error 
% \errorcontextlines
%\def\mydebug{1}
\begin{document}
\makeatletter
\grechangedim{interwordspacetext}{\unexpanded{\the\fontdimen2\font plus \the\gluestretch\gre@space@skip@interwordspacenotes minus \the\glueshrink\gre@space@skip@interwordspacenotes}}{fixed}%
\makeatother
\pagestyle{empty}
%\pagestyle{plain}	
\large
\begin{spacing}{0.7}
\begin{center}
	\section*{In Parasceve - lectio 4}
	%\small\textit{Fontevraud Gradual\\}
	\small\textit{Laudes festivae (1940) p.78}
	\vspace{5mm}
\end{center}	
\end{spacing}
%\greannotation{} 
\gresetinitiallines{1}\greannotation{Lectio 4}\grecommentary[2mm]{Ex tractátu sancti Augustíni Epíscopi super Psalmos - in Psalm 63 ad versum 2}\gabcsnippet{(f3)Ex(e) tra(f)ctá(h)tu(h) san(h)cti(h) Au(h)gu(h)stí(h)ni(h) E(h)pí(i)sco(f)pi(f) su(e)per(f) Psal(g)mos.(f) (::Z) Pro(c)te(e)xí(fg)sti(f) me,(f) De(fh)us,(h) (;) a(h) con(h)vén(h)tu(h) ma(h)li(h)gnán(h)ti(h)um,(h) a(h) mul(h)ti(h)tú(h)di(h)ne(h) o(h)pe(h)rán(h)ti(h)um(h) i(h)ni(h)qui(h)tá(i)tem.(f) (;) Iam(h) i(h)psum(h) ca(h)put(h) no(h)strum(h) in(h)tu(e)e(f)á(g)mur.(f) (:) Mul(e)ti(f) Már(h)ty(h)res(h) tá(h)li(g)a(f) pas(g)si(g) sunt,(gh) (:) sed(h) ni(h)hil(h) sic(h) e(h)lú(h)cet,(h) quó(h)mo(h)do(h) ca(g)put(f) Már(gh)ty(g)rum:(g) (;) i(e)bi(gh) mé(f)li(f)us(f) in(f)tu(f)é(f)mur,(f) (,) quod(f) il(f)li(e) ex(f)pér(g)ti(f) sunt.(f) (:) Pro(c)té(e)ctus(f) est(f) a(f) mul(f)ti(f)tú(f)di(f)ne(f) ma(f)li(e)gnán(f)ti(f)um,(f) (,) pro(f)te(f)gén(f)te(g) se(h) De(i)o,(h) (;) pro(h)te(h)gén(h)te(h) car(h)nem(h) su(h)am(h) i(h)pso(h) Fí(h)li(h)o,(h) et(h) hó(h)mi(h)ne,(h) quem(g) ge(f)ré(g)bat:(gh) (;) qui(h)a(h) fí(h)li(h)us(h) hó(h)mi(h)nis(h) est,(h) et(h) Fí(h)li(e)us(f) De(g)i(f) est.(f) (:) Fí(e)li(f)us(h) De(h)i,(h) prop(h)ter(h) for(h)mam(h) De(i)i:(f) (;) fí(h)li(h)us(h) hó(h)mi(h)nis,(h) prop(h)ter(h) for(g)mam(f) ser(gh)vi,(g) (;) ha(e)bens(gh) in(f) po(f)te(f)stá(f)te(f) pó(f)ne(f)re(f) á(f)ni(g)mam(h) su(i)am,(h) (,) et(h) re(h)cí(h)pe(e)re(f) e(g)am.(f) (:) Quid(c) e(e)i(f) po(f)tu(f)é(f)runt(e) fá(f)ce(f)re(f) i(g)ni(h)mí(i)ci?(i) (:) Oc(h)ci(h)dé(g)runt(f) cor(g)pus,(g) (;) á(f)ni(f)mam(f) non(f) oc(e)ci(f)dé(h)runt.(f) (:) In(h)tén(h)di(ef)te.(f) (:) Pa(c)rum(e) er(f)go(e) e(fh)rat,(h) (,) Dó(h)mi(h)num(h) hor(h)tá(h)ri(h) Már(h)ty(g)res(f) ver(g)bo,(gh) (;) ni(h)si(h) fir(h)má(h)ret(e) e(f)xém(g)plo.(f) (::)}\minispace%
\tlumacz{Obroniłeś mnie, Boże, od spisku złośliwych, od mnóstwa czyniących nieprawość. Tu wejrzyjmy na głowę naszą — Chrystusa. Wielu męczenników podobnie umęczonymi zostało: lecz żaden z nich nie zajaśniał tak, jak głowa męczenników patrząc na Jego cierpienia lepiej rozumiemy, co oni przenieśli. Chrystus obroniony był od mnóstwa ludzi złych; bronił Go Bóg; bronił ciała swego Syn Boży, który się stał człowiekiem, gdyż jest i Synem człowieczym i Synem Bożym. Synem Bożym z powodu natury Boskiej; synem człowieczym z powodu natury sługi, którą wziął na siebie, mając moc wydać życie swoje lub zachować je. Cóż mogli mu uczynić nieprzyjaciele? Zabili ciało, ale duszy zabić nie mogli. Rozważcie. Mało byłoby Panu zagrzewać męczenników słowem, nie dając im przykładu.}
\end{document}